\documentclass[12pt]{article}

\usepackage[a4paper,margin=2.5cm]{geometry}
\usepackage[T1]{fontenc}
\usepackage[utf8]{inputenc}
\usepackage{lmodern}
\usepackage{microtype}
\usepackage{setspace}
\usepackage{caption}
\usepackage{chngcntr}
\DeclareCaptionLabelSeparator{labdash}{ -- }
\captionsetup[figure]{labelsep=labdash}
\usepackage{booktabs}
\usepackage{array}
\usepackage{xcolor}
\usepackage{listings}
\usepackage[hidelinks]{hyperref}

\definecolor{codegray}{RGB}{245,245,245}
\setstretch{1.12}
\setlength{\parindent}{0pt}
\setlength{\parskip}{0.55em}
\renewcommand{\arraystretch}{1.2}
\lstset{basicstyle=\ttfamily\small,backgroundcolor=\color{codegray},frame=single,
  breaklines=true,showstringspaces=false,columns=fullflexible,keepspaces=true}
\newcommand{\cmd}[1]{\texttt{\detokenize{#1}}}

\counterwithin{figure}{section}
\begin{document}
\pagenumbering{gobble}
\begin{titlepage}
\centering
{\large Ministerul Educației și Cercetării al Republicii Moldova\par}
\vspace{0.45cm}
{\large Universitatea Tehnică a Moldovei\par}
\vspace{0.45cm}
{\large Facultatea Calculatoare, Informatică și Microelectronică\par}
\vspace{0.45cm}
{\large Departamentul Ingineria Software și Automatică\par}
\vspace{3cm}
{\Large\bfseries Report on Laboratory Work No. 5\par}
\vspace{0.55cm}
{\LARGE\bfseries The File System with Your Own Hands\par}
\vspace{0.35cm}
{\large Operating Systems\par}
\vfill
\begin{flushleft}
\textbf{Done by:} st. gr. FAF-241, Caț Chiril\\
\vspace{1cm}
\textbf{Checked by:}\\
asist. univ. Reițman Patricia
\end{flushleft}
\vfill
{Chișinău -- 2026\par}
\end{titlepage}

\pagenumbering{roman}
\tableofcontents
\clearpage
\pagenumbering{arabic}

\section*{Purpose and environment}
\addcontentsline{toc}{section}{Purpose and environment}
The goal was to inspect Linux path types, inode and link behavior, access permissions,
disk layout, and kernel information exposed through \cmd{/proc}. I ran the experiments
on 12 October 2026 as the normal user \cmd{class} in the installed Ubuntu guest under
UTM, not as root. The commands ran in a new directory beneath \cmd{/home/class/os-lab5}.
Their full output is preserved verbatim in the accompanying \cmd{lab5-results.txt}. The shorter command blocks below select and align relevant lines from that transcript. The CPU, RAM and disk figures describe this virtual machine.

\section{The directory tree and file types}
\begin{lstlisting}
$ ls -ld /etc /etc/passwd /bin /dev/null /tmp
lrwxrwxrwx 1 root root 7 Apr 20 11:46 /bin -> usr/bin
crw-rw-rw- 1 root root 1, 3 Oct 12 16:11 /dev/null
drwxr-xr-x 142 root root 12288 Oct 12 16:11 /etc
-rw-r--r-- 1 root root 2769 Oct  5 15:58 /etc/passwd
drwxrwxrwt 16 root root 380 Oct 12 16:15 /tmp
$ ls -l /dev/vda
brw-rw---- 1 root disk 253, 0 Oct 12 16:11 /dev/vda
$ du -sh /var/log /usr
33M   /var/log
5.7G  /usr
\end{lstlisting}
\captionof{figure}{Recorded output for path types and disk usage}

The leading \cmd{l} marks \cmd{/bin} as a symbolic link to \cmd{usr/bin}, so it
resolves to \cmd{/usr/bin}. The leading \cmd{c} marks \cmd{/dev/null} as a character
device, while \cmd{/dev/vda} starts with \cmd{b} and is a block device. The leading
\cmd{d} marks \cmd{/etc} and \cmd{/tmp} as directories; the leading \cmd{-} marks
\cmd{/etc/passwd} as a regular file. The \cmd{file} command also reported
\cmd{/etc/passwd} as ASCII text and \cmd{/bin} as a symbolic link.
\cmd{/usr} occupies about 5.7~GiB versus 33~MiB for \cmd{/var/log}: roughly
5.7~GiB more, or around 177 times as much. This reflects installed programs and
libraries in \cmd{/usr} compared with current logs in \cmd{/var/log}.

\section{Inodes, links, and allocated space}
\begin{lstlisting}
$ echo hello > notes.txt; stat notes.txt
Size: 6   Blocks: 8   IO Block: 4096   regular file
Inode: 1056330   Links: 1
$ ln notes.txt hard.txt; ls -li notes.txt hard.txt
1056330 -rw-rw-r-- 2 class class 6 Oct 12 16:15 hard.txt
1056330 -rw-rw-r-- 2 class class 6 Oct 12 16:15 notes.txt
$ rm notes.txt; cat hard.txt; ls -li hard.txt
hello
1056330 -rw-rw-r-- 1 class class 6 Oct 12 16:15 hard.txt
$ ln -s hard.txt soft.txt; ls -li hard.txt soft.txt
1056330 -rw-rw-r-- 1 class class 6 Oct 12 16:15 hard.txt
1061575 lrwxrwxrwx 1 class class 8 Oct 12 16:15 soft.txt -> hard.txt
$ rm hard.txt; cat soft.txt
cat: soft.txt: No such file or directory
$ stat soft.txt
Size: 8   Blocks: 0   symbolic link
\end{lstlisting}
\captionof{figure}{Recorded inode and link behavior}

\cmd{Inode: 1056330} identifies the file's inode. \cmd{Links: 1} is its hard-link
count: the number of directory entries naming that inode. Creating \cmd{hard.txt}
raised the count to two and gave both names the same inode. \cmd{rm notes.txt}
removed one directory entry, not the file data; \cmd{hard.txt} still named the
inode and read \cmd{hello}. The symbolic link had its own inode, 1061575. Its
contents were the eight bytes of the relative path \cmd{hard.txt}. After that
target name was removed, the symlink remained but its path resolved nowhere,
so \cmd{cat soft.txt} failed.

\begin{lstlisting}
$ mkdir d; ls -ld d
drwxrwxr-x 2 class class 4096 Oct 12 16:15 d
$ mkdir d/x d/y; ls -ld d
drwxrwxr-x 4 class class 4096 Oct 12 16:15 d
$ echo x > one.txt; ls -l one.txt; du -h one.txt
-rw-rw-r-- 1 class class 2 Oct 12 16:15 one.txt
4.0K  one.txt
$ truncate -s 1G big.img; ls -lh big.img; du -h big.img
-rw-rw-r-- 1 class class 1.0G Oct 12 16:15 big.img
0  big.img
\end{lstlisting}
\captionof{figure}{Recorded directory links and sparse file allocation}

The new directory's count of two comes from its entry in the parent and its
own \cmd{.} entry. Each new child directory has a \cmd{..} entry that points
back to \cmd{d}, raising \cmd{d}'s count to four. The two-byte file consumes
one 4~KiB filesystem block. \cmd{truncate} made a sparse 1~GiB file: it set
the logical length without writing blocks into the hole, so \cmd{du} reported
zero allocated space. The test image was removed after measurement.

\section{Permissions}
\begin{center}
\begin{tabular}{lll}
\toprule
Mode & Permission string & Typical use \\
\midrule
600 & \cmd{rw-------} & Private file readable and writable by its owner \\
644 & \cmd{rw-r--r--} & Publicly readable document, owner writable \\
755 & \cmd{rwxr-xr-x} & Executable program or traversable directory \\
700 & \cmd{rwx------} & Private directory accessible only to its owner \\
\bottomrule
\end{tabular}
\end{center}

\begin{lstlisting}
$ chmod 600 s.txt; ls -l s.txt
-rw------- 1 class class 7 Oct 12 16:15 s.txt
$ chmod 000 s.txt; cat s.txt
cat: s.txt: Permission denied
$ chmod u+r s.txt; cat s.txt
secret
$ umask; ls -ld new.txt newdir
0002
-rw-rw-r-- 1 class class 0 Oct 12 16:15 new.txt
drwxrwxr-x 2 class class 4096 Oct 12 16:15 newdir
\end{lstlisting}
\captionof{figure}{Recorded file permissions and umask}

Ownership determines which permission class applies; it does not itself grant
access. With mode 000, the owner's read bit was absent, so the normal user
could not read \cmd{s.txt}. Restoring \cmd{u+r} made it readable. The umask
\cmd{0002} removes the other-write bit from the default creation modes:
\cmd{0666 & ~0002 = 0664} for the new file and
\cmd{0777 & ~0002 = 0775} for the new directory.

\begin{lstlisting}
$ chmod u=r,go= box; ls box; cat box/f
f
cat: box/f: Permission denied
$ chmod u=x,go= box; ls box; cat box/f
ls: cannot open directory 'box': Permission denied
hi
$ ls -l /usr/bin/passwd; ls -ld /tmp
-rwsr-xr-x 1 root root 142552 Feb  2  2026 /usr/bin/passwd
drwxrwxrwt 16 root root 380 Oct 12 16:15 /tmp
$ cat /etc/shadow
cat: /etc/shadow: Permission denied
$ ls -l /etc/shadow
-rw-r----- 1 root shadow 1311 Oct  5 16:10 /etc/shadow
\end{lstlisting}
\captionof{figure}{Recorded directory traversal and special permissions}

For a directory, read permission allows listing entry names, whereas execute
permission allows searching or traversing it by a known name. Thus read-only
\cmd{box} listed \cmd{f} but blocked \cmd{box/f}; execute-only \cmd{box}
could not be listed, but \cmd{cat box/f} succeeded because the file itself
was readable. I restored \cmd{box} to mode 755 afterward.

The \cmd{s} in \cmd{passwd}'s owner-execute position is the set-user-ID bit:
the program executes with root's effective user ID. It needs that privilege to
update password data while ordinary users cannot read \cmd{/etc/shadow}.
The \cmd{t} in \cmd{/tmp} is the sticky bit. Because \cmd{/tmp} is shared and
world-writable, it restricts deletion and renaming of entries to their owner,
the directory owner, or root.

\section{Disks and the kernel as files}
\begin{lstlisting}
$ lsblk -o NAME,SIZE,TYPE,FSTYPE,MOUNTPOINTS
vda       64G disk
|-vda1     1G part vfat     /boot/efi
`-vda2  62.9G part ext4     /
$ df -h /
Filesystem  Size  Used  Avail  Use%  Mounted on
/dev/vda2    62G   12G    48G   20%  /
$ findmnt -no SOURCE,FSTYPE,TARGET /
/dev/vda2 ext4 /
$ grep -c '^processor' /proc/cpuinfo
6
$ grep '^MemTotal:' /proc/meminfo
MemTotal:        3461284 kB
\end{lstlisting}
\captionof{figure}{Recorded disk and memory observations}

The virtual disk is \cmd{/dev/vda} (64~GiB). Its first partition,
\cmd{/dev/vda1}, is a 1~GiB VFAT EFI partition mounted at \cmd{/boot/efi}.
The 62.9~GiB second partition, \cmd{/dev/vda2}, contains the ext4 root
filesystem. The root filesystem is 20\% full, with 12~GiB used and 48~GiB
available according to \cmd{df}. The VM exposes six logical processors and
3,461,284~kB of RAM, approximately 3.30~GiB.

\begin{lstlisting}
$ head -5 /proc/self/status
Name:   head
Umask:  0002
State:  R (running)
Tgid:   4301
Ngid:   0
$ ls -l /proc/self/fd
0 -> /dev/pts/0
1 -> pipe:[30151]
2 -> pipe:[30151]
3 -> /proc/4302/fd
\end{lstlisting}
\captionof{figure}{Recorded process status and file descriptors}

\cmd{/proc/self} refers to the process that opens it. Here \cmd{head}, rather
than the shell, opened \cmd{/proc/self/status}, so \cmd{Name} was \cmd{head}.
Similarly, \cmd{ls -l /proc/self/fd} listed the file descriptors of that
\cmd{ls} process: standard input was the terminal, standard output and error
were the transcript pipe, and descriptor 3 was its open fd directory.

\section*{Conclusion}
\addcontentsline{toc}{section}{Conclusion}
The experiments showed that path names and file data are separate: hard links
share an inode, while symbolic links store a path that can become dangling.
Logical file length can differ greatly from allocated disk space. Permission
checks depend on mode bits and directory traversal rights, with setuid and
sticky bits serving specific shared-system needs. Finally, disk mounts and
\cmd{/proc} exposed the guest's storage and runtime state through ordinary
filesystem interfaces.

\end{document}
